% =====================================================================
\chapter{Discussion, Limitations and Conclusion}\label{ch:conclusion}
% =====================================================================

\section{What the evidence supports}
\begin{keybox}[title={Our claim, stated precisely}]
The published scheme does not achieve infrastructure independence, escrow resilience, non-repudiation, forward secrecy or secure revocation, and its reported
costs omit the dominant work. These properties can be recovered \emph{together}, using established primitives, with lower total computation and linear
communication. The same offline trust root then extends to vehicles and to multi-TUAV handover. We do not claim new cryptographic primitives.
\end{keybox}
Each part of this claim rests on two independent kinds of evidence:
\begin{center}\small
\begin{tabularx}{\linewidth}{@{}>{\raggedright\arraybackslash}X l l@{}}
\toprule
Claim & First evidence & Second evidence\\
\midrule
Base scheme needs the TA online & protocol inspection & availability model (E1)\\
TA impersonates base-scheme UAVs & code (E7) & ProVerif Q16\\
Revoked UAV keeps the key (base) & code (E7, test) & algebraic proof (Lemma~\ref{lem:leak})\\
Our TA cannot sign under registered keys & code (test) & ProVerif Q6, Theorem~\ref{thm:escrow}\\
Our revocation locks leavers out & code (E7, test) & ProVerif Q11\\
Ours is cheaper overall & measured (E2) & analytical (Chapter~\ref{sec:perf})\\
Bounded behaviour under poisoning & measured (E3) & Proposition~\ref{prop:dos}\\
\bottomrule
\end{tabularx}
\end{center}

\section{Limitations}
\begin{enumerate}[label=\textbf{L\arabic*.},leftmargin=*]
  \item \textbf{Key replacement by a malicious TA} (Girault level~2). The TA cannot sign under a registered key, but it can issue a second key for a pseudonym or enrol a
        phantom member (ProVerif Q7, Q9). This is inherent to certificateless cryptography. It is detectable, and a manufacturer-certified device key would
        remove it.
  \item \textbf{Captured TUAV.} With the TA offline, the TUAV is the most trusted field entity. A captured TUAV learns its cell's group and pairwise keys and
        its link keys until it is evicted. It cannot forge members or learn their secret values. Threshold admission across TUAVs is future work.
  \item \textbf{Linkability within an epoch} is the price of effective local eviction; the epoch length tunes it.
  \item \textbf{Pseudonym pool.} Each epoch credential takes about 184 bytes; 288 epochs (three days of 15-minute epochs) take about 53\,KB and must be loaded before deployment.
  \item \textbf{Deferred global revocation} needs the TA to be reachable at some point; local eviction covers the gap until the epoch ends.
  \item \textbf{Larger join request} than the paper reports (225\,B versus 137\,B under our encoding; 248\,B versus 104\,B under the paper's), because ours
        carries $X$, $Y$ and an ephemeral $E$.
  \item \textbf{Measurement platform.} Timings come from a pure-Python library on a laptop whose speed varied during the runs. We therefore report ratios measured
        back to back and cross-check them with operation counts (Chapter~\ref{ch:eval}). The absolute values are not those of an optimised C implementation.
  \item \textbf{Models.} Network behaviour (link losses, mobility) is simulated with simple stochastic models rather than a radio simulator. The ProVerif
        results are symbolic, and the reductions are proof sketches.
  \item \textbf{Vehicle routine beacons} protected by the shared $\VK_a$ can be forged by any vehicle holding it; events that matter are signed.
\end{enumerate}

\section{Conclusion}
The base paper addresses a real problem, secure vehicular networking after a disaster has destroyed the roadside infrastructure, and its TUAV
architecture is a good answer. Its protocol, however, places the cloud TA inside every operation. We showed, through analysis, code and formal verification,
that this makes authentication stop when the satellite link fails, lets the TA impersonate every UAV, lets a revoked UAV keep receiving group keys, and hides
the dominant cost.

\emph{Cutting the Cord} keeps the architecture and moves the trust root offline. Escrow-free certificateless credentials let a TUAV verify UAVs and
vehicles itself, in a sound batch with a bounded fallback. Per-member ECDH keys make group-key delivery linear, forward-secure and revocation-safe.
Encrypted single-use tickets hand members between TUAVs without the TA. ProVerif verifies the intended properties and pins down the one residual TA power.
On the same curve and library, our TUAV authenticates and keys a swarm one to two orders of magnitude faster than the base scheme's TA-inclusive cost,
with no satellite traffic and kilobytes instead of megabytes of key traffic.

\section{Future work}
\begin{itemize}
  \item Threshold trust across TUAVs, so that one captured TUAV cannot admit members alone.
  \item Level-3 trust: bind each RG1 record to a manufacturer-certified device key.
  \item A constant-time implementation on a 128-bit curve, measured on embedded UAV hardware.
  \item Integration with a radio and mobility simulator (for example ns-3) and a field trial with a tethered drone.
  \item Post-quantum credentials (lattice-based signatures) for long-lived deployments.
\end{itemize}
